\documentclass[11pt]{article}

\usepackage[margin=25mm]{geometry}

\usepackage{amsmath,amssymb,fontspec,unicode-math}

\IfFontExistsTF{DejaVu Serif}{\setmainfont{DejaVu Serif}}{\setmainfont{Latin Modern Roman}}

\IfFontExistsTF{Latin Modern Math}{\setmathfont{Latin Modern Math}}{\setmathfont{DejaVu Math TeX Gyre}}

\usepackage{xurl,hyperref}

\setlength{\parindent}{0pt}\setlength{\parskip}{6pt}

\begin{document}

\title{DSAS: A Universal Plug-and-Play Framework for Attention Optimization in Multi-Document Question Answering}

\author{ELVIA preliminary AI analysis}\date{21 October 2026}\maketitle

Preliminary assessment: Minor.\par

Independent scientific review remains pending. Author adoption has not been requested.

Signed-logit scaling does not guarantee probability suppression; the current post-publication adapter also contains mask/index defects. Dense complexity, exported analysis groups and partial-resume accounting require correction. Published model improvements are not reproduced or refuted.

\section*{Review scope}

Claim-focused reading of the original paper and relevant author-released passages; direct mathematical or static-source reasoning. No author program, earlier verifier, model or benchmark was executed during this check.

\section*{Problem 1: Scaling a negative logit increases its attention probability}

Source: Eqs.(3),(13)–(14), PDF pp.4,6

Assessment: supported; impact: minor

\subsection*{Source claim (paraphrase)}

CGW/RAS multiply selected raw scores by \(w\in[0.7,1]\) before softmax, with a suppression interpretation.

\subsection*{Analysis}

For legal causal logits \((-2,0)\), scaling the first by \(0.7\) changes its probability from \(e^{-2}/(1+e^{-2})=0.1192029220\) to \(e^{-1.4}/(1+e^{-1.4})=0.1978161114\). Relative odds grow by \(e^{0.6}\). Arbitrary finite scores admit a query/key factorization, so positivity is not a hidden dot-product restriction. Positive logits can instead be suppressed, and empirical average improvement remains possible.

\subsection*{Finding and scope}

The operation is not a general probability-suppression guarantee; no actual model-quality reversal is established.

\subsection*{Proposed correction}

If probability gating is intended, additive log gates \(z_j+\log w_j\) give odds scaled by \(w_j/w_k\). A common gate below one reduces a selected group’s total mass against its unchanged complement. Arbitrary heterogeneous gates need more careful coordinate claims, and this changed intervention requires experiments.

A proposed correction is retained where stated; independent verification and author adoption remain unestablished.

\section*{Problem 2: Chained advanced indexing leaves the intended masks unchanged}

Source: adapter\_for\_train.py, lines 89–91,104–113; Eqs.(13)–(14), PDF p.6

Assessment: supported; impact: minor

\subsection*{Source claim (paraphrase)}

The assignments index paragraph/question ranges in two successive operations, then apply an initially false mask.

\subsection*{Analysis}

Advanced range indexing returns a copy: assigning the second-indexed temporary cannot set the original mask. Four paragraphs of \(10\) tokens, question interval \([40,44)\), target \(44\), and tensor shape \((1,1,45,45)\) meet top-\(10\) requirements. A causal paragraph block has \(100\) intended mask entries but none are set. The separate direct target assignment still acts, so this is not a wholly bypassed adapter. An additional source threshold \((\beta+1)/2=0.85\), absent from Eq.(13), skips gates above that value; an index-only repair does not fix it.

\subsection*{Finding and scope}

Current question-to-paragraph CGW and RAS branches do not perform the advertised updates. No conclusion about the 2025 ablation code follows.

\subsection*{Proposed correction}

For contiguous ranges, assign both dimensions in one basic slice, and separately remove or justify the extra threshold if Eq.(13) is intended. Index-only patches retain the signed-logit issue.

A proposed correction is retained where stated; independent verification and author adoption remain unestablished.

\section*{Problem 3: Rank slots are confused with paragraph identifiers}

Source: adapter\_for\_train.py, lines 59–76; Eqs.(8)–(11), PDF pp.5–6

Assessment: supported; impact: minor

\subsection*{Source claim (paraphrase)}

The loop computes a positional exponent for paragraph indices[i], writes it to gama[i], then consumes gama at the paragraph ID.

\subsection*{Analysis}

With descending paragraph IDs \((3,2,1,0)\), intended exponents are \[(1,1,0.3808301389,0.2391125411)\], but the assignment produces \[(0.2391125411,0.3808301389,1,1)\]. The later read at IDs \(3,2\) uses \(1\) for both. Under the source convention this changes gates from approximately \((0.7,0.72282,0.79275,1)\) to \((0.7,0.78730,0.89866,1)\), and changes the above-mean key set. IDs equal to rank slots hide the discrepancy.

\subsection*{Finding and scope}

The paragraph-specific operation is misaligned. These controlled source-formula values are not measurements of a trained model.

\subsection*{Proposed correction}

Write the computed exponent at its paragraph ID. This fixes alignment only; mask, threshold and signed-score issues remain separate.

A proposed correction is retained where stated; independent verification and author adoption remain unestablished.

\section*{Problem 4: Dense RAS is omitted from complexity accounting}

Source: Appendix E/Table 6, PDF p.25; adapter\_for\_train.py, lines 42,89,104,111

Assessment: supported; impact: minor

\subsection*{Source claim (paraphrase)}

The whole score refinement is described as \(O(LQ)\) time and \(O(C)\) space.

\subsection*{Analysis}

Fix question length \(Q\) and paragraph count \(C\), and increase paragraph lengths. A key/irrelevant causal rectangle in Eq.(14) contains \(\Theta(L^2)\) scores, requiring that many scalar writes when materialized densely. The source additionally allocates full boolean masks and an \(L\times L\) head sum: one mask alone has \(L^2\) elements even when its advanced-index assignment fails. A persistent paragraph vector can use \(O(C)\), but does not account for these temporaries.

\subsection*{Finding and scope}

The dense-stage time/workspace description omits quadratic terms. Fused attention could change incremental memory, but no such implementation is established here; no wall-clock slowdown is inferred.

\subsection*{Proposed correction}

Separate gate extraction/storage from score application. Include worst-case dense quadratic work and current full-mask temporary storage.

A proposed correction is retained where stated; independent verification and author adoption remain unestablished.

\section*{Problem 5: The exported BAD group includes partly correct answers}

Source: Appendix A.4, PDF p.23; Section 2.3, p.4; evaluation.py, lines 103–108

Assessment: supported; impact: minor

\subsection*{Source claim (paraphrase)}

The appendix defines GOOD by token F1 \(1\) and BAD by precision \(0\); the source groups exact-string matches as GOOD and all other answers as BAD.

\subsection*{Analysis}

Against reference “New York”, prediction “York” has precision \(1\), recall \(1/2\), F1 \(2/3\): it belongs to neither appendix extreme but is exported BAD. “York New” has token F1 \(1\) but a different string order and is also exported BAD. Main-text exact-match wording can justify a stricter GOOD rule, but cannot justify calling a partially overlapping answer no-overlap BAD. Aggregate numerical F1 need not change.

\subsection*{Finding and scope}

Exported group composition differs from the stated analysis. The connection to historical information-flow plots remains unestablished.

\subsection*{Proposed correction}

Choose and state the GOOD rule, and use an explicit zero-precision condition for BAD, leaving partial cases outside the extreme groups.

A proposed correction is retained where stated; independent verification and author adoption remain unestablished.

\section*{Problem 6: Partial resume truncates earlier predictions}

Source: method/DSAS.py, lines 69–101

Assessment: supported; impact: minor

\subsection*{Source claim (paraphrase)}

The script reads existing IDs and excludes them from remaining work, then opens the same output in write mode.

\subsection*{Analysis}

For planned IDs A,B with A already saved, the remaining set is B. Opening with write truncates A; after B is written the file contains only B. A complete existing run exits earlier and an uninterrupted new run has no old records to lose.

\subsection*{Finding and scope}

The conditional I/O path loses predictions and can change a later evaluation denominator. Historical interrupted use is not established.

\subsection*{Proposed correction}

Append new records after excluding old IDs, or atomically write a complete merged set. The neighboring LongBench append branch provides a local control.

A proposed correction is retained where stated; independent verification and author adoption remain unestablished.

\section*{Limits}

Preliminary AI assessment; independent human scientific adjudication remains pending.\par

Analytic witnesses and static release checks do not remeasure the published benchmark results.\par

The inspected source commit is dated 31 August 2026, after publication; its historical use for the 2025 tables is unestablished.\par

\section*{Source identity}

NeurIPS 2025 official proceedings PDF (retained local copy)

\url{https://proceedings.neurips.cc/paper_files/paper/2025/file/ff1789b3eb6cbc2ee6a9bdd9393b9ed2-Paper-Conference.pdf}

Reviewed PDF SHA-256: \nolinkurl{eebafa5029d42e3a9b5c9a74b9a0df1060b278342bfc10906d3b4e6222ba422b}

The SHA-256 identifies the exact retained PDF inspected in this rapid source/argument check. Original supplied report bytes are retained locally. No unambiguous original PDF retrieval timestamp was recovered; file modification time was not substituted.

\end{document}
